\documentclass[a4paper,10pt,twoside]{article}
\usepackage[margin=0.75in, top=1.2in, bottom=1.4in, headheight=15pt, headsep=0.4in, footskip=0.6in]{geometry} % Adjusted margins and spacing for header and footer
% === DYNAMIC PREAMBLE ===
\usepackage{polyglossia}
\setdefaultlanguage{english}

        \usepackage{fontspec}
        \setmainfont{NotoSans}[
            Path=../../../Fonts/,
            UprightFont    = *-Regular ,
            BoldFont       = *-Regular ,
            ItalicFont     = *-Regular ,
            BoldItalicFont = *-Regular ,
            BoldFeatures   = {FakeBold=1.5} ,
            ItalicFeatures = {FakeSlant=0.2} ,
            BoldItalicFeatures = {FakeBold=1.5, FakeSlant=0.2}
        ]
% === END DYNAMIC PREAMBLE ===
\usepackage[dvipsnames,svgnames,x11names]{xcolor}
\usepackage{amsmath,amssymb,mathtools}
\usepackage{tikz}
\usepackage{booktabs}
\usepackage{graphicx}
\usepackage{hyperref}
\usepackage{fancyhdr}
\usepackage{titlesec}
\usepackage{multicol}

\definecolor{ImageNavy}{HTML}{001f5c}

% === LINK COLORS ===
\hypersetup{
    colorlinks=true,
    linkcolor=black,
    urlcolor=black,
    citecolor=black
}

% ==========================================
% FONT COMMANDS (XeLaTeX / fontspec compatible)
% ==========================================
\newcommand{\journalmetafont}{\fontsize{10}{12}\franklingothicbookfont\selectfont}
\newfontfamily\timesfont{Times New Roman}[Path=../../../Fonts/, UprightFont=*, BoldFont=*-Bold, ItalicFont=*-Italic, BoldItalicFont=*-BoldItalic]
\newfontfamily\calibrifont{Calibri}[Path=../../../Fonts/, UprightFont=*-Regular, BoldFont=*-Bold, ItalicFont=*-Italic, BoldItalicFont=*-BoldItalic]
\newfontfamily\franklingothicfont{FRADM}[Path=../../../Fonts/, UprightFont=*, ItalicFont=*IT]
\newfontfamily\franklingothicbookfont{FRABK}[Path=../../../Fonts/, UprightFont=*, ItalicFont=*IT]
\newfontfamily\arialfont{arial}[Path=../../../Fonts/, UprightFont=arial, BoldFont=arialbd, ItalicFont=ariali, BoldItalicFont=arialbi, Extension=.ttf]
\newcommand{\articletypefont}{\fontsize{9}{13}\timesfont\itshape\bfseries\selectfont}
\newcommand{\articleTitleFont}{\fontsize{17}{20}\franklingothicfont\bfseries\color{ImageNavy}\selectfont}
\newcommand{\authorfont}{\fontsize{11}{15}\timesfont\bfseries\selectfont}
\newcommand{\affiliationfont}{\fontsize{10}{13}\arialfont\selectfont}
\newcommand{\emailfont}{\fontsize{9}{11}\arialfont\selectfont}
\newcommand{\headingfont}{\fontsize{16}{14}\calibrifont\bfseries\selectfont}
\newcommand{\subheadingfont}{\fontsize{10}{14}\bfseries\selectfont}
\newcommand{\contentfont}{\fontsize{10}{13}\calibrifont\selectfont}
\newcommand{\footerfont}{\fontsize{8}{11}\selectfont}
\newcommand{\citationfont}{\fontsize{8}{11}\selectfont}

% ==========================================
% SECTION FORMATTING
% ==========================================
\titleformat{\section}{\headingfont}{\thesection}{1em}{}
\titleformat{\subsection}{\subheadingfont}{\thesubsection}{1em}{}
\titlespacing{\section}{0pt}{6pt}{8pt}

% ==========================================
% EXTRACT AUTHOR PARTS FOR CITATION
% ==========================================

% ==========================================
% HEADER & FOOTER --- DEFAULT (subsequent pages)
% ==========================================
\pagestyle{fancy}
\fancyhf{}
\renewcommand{\headrulewidth}{0pt}
\renewcommand{\footrulewidth}{0.4pt}
\fancyfoot[L]{%
    \citationfont
    \vspace{0.1cm}
    \noindent \textbf{Citation:} Smith J (17). Machine Learning: Transforming Science, Business, and Life Transforming Science, Business, and Life. J. AI Res.. 06: 56.
}
\fancyhead[LE]{\fontsize{9}{11}\timesfont\selectfont \thepage\quad J. AI Res.}
\fancyhead[RO]{\fontsize{9}{11}\timesfont\selectfont John Smith\quad \thepage}

% ==========================================
% FIRST PAGE STYLE
% ==========================================
\fancypagestyle{firstpage}{
    \fancyhf{}
    \renewcommand{\headrulewidth}{0pt}
    \renewcommand{\footrulewidth}{0pt}
    \fancyfoot[C]{
        \begin{minipage}{\textwidth}
            \hrule height 0.4pt
            \vskip 0.15cm
            \noindent \footerfont \textbf{Received:} 18-Jul-2026, Manuscript No. 2024-001; \textbf{Editor assigned:} 20-Jul-2026, PreQC No. 2024-001 (PQ); \textbf{Reviewed:} 03-Aug-2026, QC No. 2024-001; \textbf{Revised:} 10-Aug-2026, Manuscript No. 2024-001 (R); \textbf{Published:} 2026-08-17
            \vskip 0.15cm
            \hrule height 0.4pt
            \vskip 0.15cm
            \noindent \citationfont \textbf{Citation:} Smith J (17). Machine Learning: Transforming Science, Business, and Life Transforming Science, Business, and Life. J. AI Res.. 06: 56.
        \end{minipage}
    }
}

% ==========================================
% MAIN DOCUMENT
% ==========================================
\begin{document}

\thispagestyle{firstpage}

\vspace*{-0.6in} % Pulls first page up to negate the global top=1.2in margin (reverts to ~0.75in effective top margin)

% --- Top header block: logo + journal metadata ---
\noindent
\begin{minipage}[t]{3.4cm}
    \vspace{0pt}
    \includegraphics[height=3.81cm, width=3.4cm]{../../../Logo/intresjournal.jpeg}
\end{minipage}%
\hspace{0cm}
\begin{minipage}[t]{\dimexpr\textwidth-3.7cm\relax}
    \vspace{0pt}
    \raggedright
    \journalmetafont
    Journal of AI Research (ISSN: 1234-567X) Vol. 06(02) pp. 56, 2026-08-17 \\
    DOI: \href{https://doi.org/10.1234/jair.2026.001}{https://doi.org/10.1234/jair.2026.001} \\
    Available online \href{https://example.com/journal-ai}{https://example.com/journal-ai} \\
    Copyright \copyright\ 17 International Research Journals

    \vspace{1.7cm}
    {\articletypefont Research Article}
\end{minipage}

\vspace{0.4cm}

% --- Title, author, affiliation block ---
\begin{center}
    {\articleTitleFont Machine Learning: Transforming Science, Business, and Life Transforming Science, Business, and Life\par}
    \vspace{0.25cm}
    {\authorfont John Smith\textsuperscript{*}\par}
    \vspace{0.15cm}
    {\affiliationfont Computer Science\par}
    {\emailfont E-mail: \href{mailto:john.smith@example.com}{john.smith@example.com}\par}
\end{center}

\vspace{0.4cm}

% ==========================================
% BODY --- TWO-COLUMN LAYOUT
% ==========================================
\contentfont
\setlength{\columnsep}{0.7cm}
\setlength{\parindent}{0pt}
\setlength{\parskip}{0.5em}
\begin{multicols}{2}

    \section*{INTRODUCTION}
    \begin{otherlanguage}{english}
    The pervasive integration of machine learning (ML) into scientific endeavors is fundamentally reshaping the landscape of research and discovery across numerous disciplines. This transformative technology is proving instrumental in accelerating the pace of innovation, refining experimental methodologies, and unlocking unprecedented capabilities in data analysis. In the realm of drug discovery, ML algorithms are adept at identifying novel therapeutic targets and predicting the efficacy of potential drug candidates, thereby significantly shortening development cycles and reducing costs [1].

Similarly, materials science is witnessing a revolution with ML's ability to predict the properties of new materials, guide synthesis processes, and optimize material performance for specific applications. This predictive power allows researchers to explore vast combinatorial spaces of potential materials that would be intractable through traditional experimentation alone [1].

In climate modeling, ML is contributing to more accurate predictions of weather patterns, climate change impacts, and the effectiveness of mitigation strategies. By analyzing complex datasets from various sources, ML models can discern intricate relationships and provide insights that enhance our understanding of global climate systems [1].

The field of natural language processing (NLP) has seen remarkable advancements due to deep learning techniques. These techniques enable machines to understand, interpret, and generate human language with increasing sophistication, driving progress in applications such as machine translation and sentiment analysis [2].

Cybersecurity is another critical domain benefiting from ML's capabilities. The ability of ML algorithms to process vast quantities of network traffic and system logs allows for the real-time detection of anomalies indicative of malicious activity, thus bolstering defenses against evolving cyber threats [3].

The healthcare sector is increasingly leveraging ML for a variety of applications, ranging from enhanced diagnostic accuracy through medical image analysis to the development of personalized treatment plans and accelerating drug development pipelines, ultimately aiming to improve patient outcomes and streamline healthcare operations [4].

Robotics and autonomous systems are being propelled forward by reinforcement learning (RL), a subfield of ML. RL enables robots to learn complex control policies through experience, allowing them to navigate dynamic environments, perform intricate tasks, and interact more effectively with humans and their surroundings [5].

Financial markets are also being influenced by ML, particularly in the areas of forecasting and algorithmic trading. ML models can analyze historical data and market sentiment to predict stock price movements and develop automated trading strategies, aiming to capitalize on market opportunities and manage risk [6].

Intelligent transportation systems (ITS) are being enhanced by ML through applications such as traffic prediction, route optimization, and autonomous vehicle control. By integrating ML with sensor data and communication networks, ITS can improve road safety, reduce congestion, and boost overall transportation efficiency [7].

In agriculture, ML is driving the adoption of precision farming techniques. By analyzing data from various sources, ML models can optimize resource allocation for irrigation, fertilization, and pest control, leading to increased crop yields, reduced environmental impact, and enhanced food security [8].
    \end{otherlanguage}

    \vspace{0.4em}

    \section*{DESCRIPTION}
    \begin{otherlanguage}{english}
    The application of machine learning (ML) in scientific research offers a paradigm shift, enabling the exploration of complex phenomena and the acceleration of discovery processes. Its ability to identify subtle patterns and relationships within large datasets allows scientists to formulate hypotheses and design experiments with greater precision and efficiency. For instance, in the pharmaceutical industry, ML accelerates drug discovery by predicting molecular interactions and identifying potential drug candidates, thereby significantly reducing the time and cost associated with traditional research methods [1].

In materials science, ML empowers researchers to design novel materials with desired properties by predicting material behavior and guiding synthesis. This computational approach complements experimental work, allowing for the rapid exploration of vast material design spaces and leading to the development of advanced materials for diverse applications [1].

Climate science benefits immensely from ML's capacity to analyze large-scale environmental data. ML models can improve the accuracy of climate predictions, identify key drivers of climate change, and assess the potential impact of various intervention strategies, contributing to more informed policy decisions [1].

Deep learning, a sophisticated subset of ML, has revolutionized natural language processing (NLP). Its advancements have enabled machines to comprehend and generate human language with unprecedented accuracy, leading to significant improvements in applications such as real-time language translation, sophisticated chatbots, and advanced text summarization tools [2].

Cybersecurity is a crucial area where ML plays a vital role in threat detection and mitigation. ML algorithms can analyze network traffic in real-time to identify anomalies that signal cyberattacks, such as malware infections or phishing attempts, thereby providing an early warning system and enabling rapid response [3].

In the healthcare domain, ML is instrumental in improving patient care through applications like disease diagnosis, personalized treatment planning, and drug development. By analyzing medical images, patient data, and genomic information, ML models can assist clinicians in making more accurate diagnoses and tailoring treatments to individual patients [4].

Reinforcement learning (RL) is a key enabler for advanced robotics and autonomous systems. RL allows robots to learn optimal behaviors through trial and error in complex and dynamic environments, facilitating their deployment in tasks requiring dexterity, navigation, and human-robot interaction [5].

Financial forecasting and algorithmic trading have been significantly impacted by ML techniques. Models are employed to predict stock market trends, analyze market sentiment, and automate trading decisions, aiming to enhance profitability and manage investment risks more effectively [6].

Intelligent transportation systems (ITS) leverage ML for sophisticated traffic management and control. Applications include predicting traffic flow, optimizing routes for reduced travel times, and enabling the safe operation of autonomous vehicles, all contributing to more efficient and safer transportation networks [7].

Agriculture is being transformed by ML through precision farming techniques. By analyzing data from sensors, drones, and satellites, ML models help optimize irrigation, fertilization, and pest management, leading to higher crop yields, reduced resource consumption, and improved sustainability in food production [8].
    \end{otherlanguage}

    \vspace{0.4em}

    \section*{CONCLUSION}
    \begin{otherlanguage}{english}
    This collection of research explores the widespread applications of machine learning across diverse fields. Machine learning is revolutionizing scientific discovery, enabling accelerated research in areas like drug discovery, materials science, and climate modeling. Deep learning advancements are transforming natural language processing, while ML enhances cybersecurity by detecting threats and improves healthcare through diagnostics and personalized treatments. In robotics, reinforcement learning drives autonomous systems, and financial markets utilize ML for forecasting and trading. Furthermore, ML optimizes intelligent transportation systems and agricultural practices, contributing to efficiency and sustainability. Emerging applications also extend to renewable energy forecasting and educational technology.
    \end{otherlanguage}

    \vspace{0.4em}

    % ==========================================
    % REFERENCES
    % ==========================================
    \section*{REFERENCES}
    {\citationfont
                \noindent\hangindent=1.5em\hangafter=1
                    Marcos VR, Sergej FGW, Alana LVJ et al.
                (2021).
                    \href{https://pubmed.ncbi.nlm.nih.gov/33753790/}{Machine learning for scientific discovery.}
                \textit{Nature},
                    591(591):523-530.
                \par\vspace{3pt}

            \noindent\hangindent=1.5em\hangafter=1
                    Quan L, Yongliang S, Wenfeng C et al.
                (2022).
                    \href{https://ieeexplore.ieee.org/abstract/document/9530071/}{A Survey of Deep Learning for Natural Language Processing.}
                \textit{IEEE Transactions on Pattern Analysis and Machine Intelligence},
                    44(12):8275-8291.
                \par\vspace{3pt}

            }

\end{multicols}

\end{document}